% Generated draft: requires mathematical and source review.
\documentclass[11pt]{article}
\usepackage[margin=25mm]{geometry}
\usepackage{fontspec}
\setmainfont{Latin Modern Roman}
\usepackage{amsmath,amssymb,mathtools,unicode-math}
\setmathfont{Latin Modern Math}
\usepackage{enumitem,needspace,xurl,hyperref}
\setlength{\parindent}{0pt}
\setlength{\parskip}{5pt}
\setlength{\emergencystretch}{3em}
\newcommand{\N}{\mathbb{N}}
\newcommand{\Z}{\mathbb{Z}}
\newcommand{\Q}{\mathbb{Q}}
\newcommand{\R}{\mathbb{R}}
\newcommand{\C}{\mathbb{C}}
\newcommand{\F}{\mathbb{F}}
\newcommand{\A}{\mathbb{A}}
\newcommand{\PP}{\mathbb{P}}
\newcommand{\E}{\mathbb{E}}
\newcommand{\eps}{\varepsilon}
\newcommand{\dd}{\,\mathrm{d}}
\begin{document}
\section*{UnsolvedMath ID: 30005248}
\textbf{Source problem code:} catalogue entry
\subsection{Definitions and mathematical statement}
The supplied record asks whether the dependence on a reference distribution \(P_0\) observed in testing rates also occurs when estimating the functional \(\rho(P,P_0)\). It does not specify a statistical experiment, distribution class, loss, localization, or admissible family of functionals. Consequently, it determines a research question rather than a uniquely specified minimax problem.

A conditional framework for making the estimation question precise is as follows. Let \(\mathcal P\) be a specified class of probability distributions, let \(P_0\in\mathcal P\), and let \(\rho:\mathcal P\times\mathcal P\to T\) take values in a specified measurable target space \(T\). For each sample-size index \(n\) and unknown \(P\in\mathcal P\), let \(Q_{P,n}\) denote the law of the observed data. Specify a loss \(L:T\times T\to[0,\infty]\) and a nonempty local parameter set \(\mathcal U_n(P_0)\subseteq\mathcal P\). Whenever the requisite measurability conditions hold, define
\[
R_n(P_0)=\inf_{\widehat\theta_n}\;\sup_{P\in\mathcal U_n(P_0)}
\mathbb E_{Q_{P,n}}\!\left[L\!\left(\widehat\theta_n,\rho(P,P_0)\right)\right],
\]
where the infimum is over admissible measurable estimators based on the observed data. Whether those estimators may use \(P_0\) must also be specified. These ingredients provide a definition template; they are not additional hypotheses established by the supplied source excerpt.

\textbf{Research question.} Determine for which functionals and distribution classes the local minimax estimation rate exhibits substantial dependence on \(P_0\). The registry formulation additionally asks how fuzzy-testing rates transition to functional-estimation rates as the null radius increases.

\textbf{Unresolved qualifications.} The source does not define substantial dependence, the relevant asymptotic regime, the local parameter sets, or the null radius. It also does not specify the fuzzy-testing null and alternative hypotheses, their separation, or the permitted testing errors. Thus the second question has no uniquely defined testing rate or transition in the supplied evidence. Its explicit formulation appears in the registry statement but not in the supplied original quotation.
\subsection{Short English statement}
The question is whether estimating \(\rho(P,P_0)\) can have a rate that depends substantially on the reference distribution \(P_0\), as the supplied account says occurs in testing. The registry also asks about a transition from fuzzy testing to estimation as the null radius grows. Both questions require further statistical specifications before they become precise minimax problems.
\subsection{Sources}
\begin{itemize}
\item Local UnsolvedMath registry, permanent ID \(30005248\), title \emph{Local Complexity of Functional Estimation}, source problem code \texttt{catalogue entry}. The supplied record lists proposed year 2022 and supplies no proposer attribution.
\item Supplied source citation: \emph{Mini-Workshop: Mathematical Foundations of Robust and Generalizable Learning} (2023), Oberwolfach Reports, pp.\ 2659--2686, DOI \texttt{10.4171/owr/2022/46}. \url{https://doi.org/10.4171/owr/2022/46}. The literature-source entry labels this same reference \emph{Oberwolfach Report 46/2022}.
\item The supplied original excerpt asks whether the rate for estimating \(\rho(P,P_0)\) depends strongly on \(P_0\). It is immediately followed by the text ``Breaking the sample complexity barrier in reinforcement learning Yuxin Chen''. The record does not establish whether this trailing text belongs to the problem or marks a subsequent contribution; it does not establish attribution of this problem to Yuxin Chen.
\item Literature assessment dated exactly \texttt{2026-08-22}: the supplied curation labels the problem open and reports that the source describes the corresponding dependence in functional estimation as unknown. This is supplied dated curation, not an independently verified literature assessment. Its listed source is \url{https://doi.org/10.4171/owr/2022/46}.
\item The registry labels its statement \texttt{corrected\_verified} and reports that it was checked against the official report. That verification claim is supplied metadata and has not been independently verified here.
\end{itemize}
\end{document}
